\sscp{Labelling of language names}{02-04-02}
Language names may be mislabelled,
misleading, or are actually the names of villages.
For example, many readers are unaware that \citet{Collins1983}
tends to label his data according to the village name in which
a word list was taken or from which the consultant came,
and not necessarily the language name.
So some readers are unaware that several data sets are from different
varieties of the same language such as \ili{Alune} or \ili{Wemale}.
Only a careful reading may help sort that out.
Even those reasonably familiar with the wider region may be easily confused.

Similarly \citet{BlustTrussel2020} list data from “Tifu”.
There is no Tifu language,
and it should not be treated as a separate witness.
The data are clearly from the \ili{Buru} language.
Tifu is a small village on the south coast of \ili{Buru} where Dutch
UZV (\it{Utrechtse Zendings Vereniging}) missionaries Hendricks and Schut resided.
They both published on the \ili{Buru} language
and culture, and were the main sources for
\citeauthor{Stresemann1927}'s (\citeyear{Stresemann1927}) \ili{Buru} data.

In several of our subgroups,
varieties that share the same language name in the literature
may have patterned differences in their sound correspondences
that are important for historical-comparative work.
This is true for \ili{Alune},
\ili{Wemale}, \ili{Kei}, “Rote”, “\ili{Meto}”,
“\ili{Manggarai}”, and other large or geographically widespread languages.

In some cases, we have been unable to confidently identify which
language a particular set of data represents or its geographic
origin due to problems in labelling.
Thus, \citet{Tauern1931} has short grammar sketches of a number
of languages of Seram.
One of these is labelled “Uhei-Kaẖlakim”
and another is labelled “\ili{Liambata}”.
Based on comparison with data in other sources,
these are probably varieties of \ili{Seti} (\srf{10-03-04}),
but we are not completely confident.
\citet{Stresemann1927} also has “\ili{Liambata}” data which is different
to that in \citet{Tauern1931},
and, to further complicate matters,
\citet{Tauern1931} also refers to his “\ili{Liambata}” as “\ili{Bonfia}”,
which is the Dutch-era label for \ili{Masiwang} [bnf].
Comparison with sources which unambiguously present \ili{Masiwang} data
(e.g. \citealp{LeCocqdArmandville1901} and
\citealp{Collins1986})
shows that the “\ili{Liambata}” data in \citet{Tauern1931}
is definitely \it{not} \ili{Masiwang},
but it is easy to see how confusion
and mis-identification can arise.

